\clearpage
\VerdierPhysicalPageAnchor{0071}%
\VerdierSemanticAnchor{ch1-2-3}
\subsection{複体の圏上の次数付き構造}

\VerdierSemanticAnchor{ch1-2-3-1}
\noindent\textbf{2.3.1.}\quad
$Y$ を $n$-重複複体とし、$i\in[n]$ とする。次のように定める：

\VerdierSemanticAnchor{ch1-2-3-1-eq1}
\begin{equation}
  \begin{aligned}
    (T_iY)^\sigma&=Y^{\sigma+e_i},
      &&\sigma\in\mathbb Z[n],\\
    d_{T_iY}^{\sigma,j}&=\nu(i,j)d_Y^{\sigma+e_i,j},
      &&j\in[n],
  \end{aligned}
  \tag{2.3.1.1}
\end{equation}
ここで、$i\ne j$ のとき $\nu(i,j)=1$、$i=j$ のとき $\nu(i,j)=-1$ とする。

\VerdierSemanticAnchor{ch1-2-3-1-morphism-shift}
$n$-重複複体の任意の射
$f\colon Y\longrightarrow Y'$ に対して、次のように定める：

\VerdierSemanticAnchor{ch1-2-3-1-eq2}
\begin{equation}
  (T_if)^\sigma=f^{\sigma+e_i},
  \qquad \sigma\in\mathbb Z[n].
  \tag{2.3.1.2}
\end{equation}

\VerdierSemanticAnchor{ch1-2-3-1-translation-functor}
これにより関手
$T_i\colon\operatorname{comp}^{n}(C)\longrightarrow
\operatorname{comp}^{n}(C)$ を定義し、\emph{方向 $i$ のシフト関手}と呼ぶ。

\VerdierSemanticAnchor{ch1-2-3-1-n-equals-one}
$n=1$ のとき、記号 $TY$、$Tf$、または $Y[1]$、$f[1]$ を用いる。
方向 $i$ のシフト関手は $\operatorname{comp}^{n}(C)$ の自己同型である。

\VerdierSemanticAnchor{ch1-2-3-1-eq3}
関手 $T_i$（$i\in[n]$）は互いに可換である：
\begin{equation}
  T_iT_j=T_jT_i,
  \qquad \forall i,j\in[n].
  \tag{2.3.1.3}
\end{equation}

\VerdierSemanticAnchor{ch1-2-3-1-action}
したがって、群 $\mathbb Z[n]$ の圏
$\operatorname{comp}^{n}(C)$ 上の作用を次のように定義できる：

\VerdierSemanticAnchor{ch1-2-3-1-eq4}
\begin{equation}
  \sigma=\sum_{i\in[n]}l_ie_i,
  \qquad
  T(\sigma)=T_1^{l_1}\cdots T_n^{l_n}.
  \tag{2.3.1.4}
\end{equation}

\VerdierSemanticAnchor{ch1-2-3-1-strict-graded-category}
以下、記号 $\operatorname{comp}^{n}(C)$ は、上述の厳密な
$\mathbb Z[n]$-圏構造を備えた圏 $\operatorname{comp}^{n}(C)$ を表す。

\VerdierSemanticAnchor{ch1-2-3-2}
\noindent\textbf{2.3.2.}\quad
$\varphi\colon[n]\longrightarrow[m]$ を写像、$Y$ を $C$ の $n$-重複複体、
$i$ を $[n]$ の元とする。話を具体的にするため、$C$ において可算直和が
表現可能であると仮定する。

\VerdierSemanticAnchor{ch1-2-3-2-equal-objects}
まず、対象
$\left(\int_\varphi^\Sigma T_iY\right)^\tau$ と
$\left(T_{\varphi(i)}\int_\varphi^\Sigma Y\right)^\tau$ が等しいことに注意する。
実際、定義により、これらはいずれも次に等しい：

\VerdierSemanticAnchor{ch1-2-3-2-direct-sum}
\begin{equation*}
  \bigoplus_{\varphi(\sigma)=\tau+e_{\varphi(i)}}Y^\sigma.
\end{equation*}
そこで、次のように定める：

\VerdierSemanticAnchor{ch1-2-3-2-eq1}
\begin{equation}
  \begin{aligned}
    \bigl(p_{i,\varphi}(Y)\bigr)^\tau\colon
    \left(\int_\varphi^\Sigma T_iY\right)^\tau
      &\longrightarrow
    \left(T_{\varphi(i)}\int_\varphi^\Sigma Y\right)^\tau,
      &&\tau\in\mathbb Z[m],\\[2pt]
    \bigl(p_{i,\varphi}(Y)\bigr)^\tau
      &=\bigoplus_{\varphi(\sigma)=\tau+e_{\varphi(i)}}
        \varepsilon(\sigma,i)\varepsilon(\varphi(\sigma),\varphi(i))
        \operatorname{id}_{Y^\sigma}.
  \end{aligned}
  \tag{2.3.2.1}
\end{equation}
